\section{Introduction}
\label{sec:intro}

Deep neural networks (DNNs) achieve strong performance across computer vision \cite{he2016deep,dosovitskiy2020image}, but this performance comes with high computational cost, large parameter counts and substantial memory demand \cite{hoefler2021sparsity}. These costs make deployment difficult on mobile, embedded and edge devices \cite{howard2017mobilenets,ma2018shufflenet}. Model compression addresses this through quantization \cite{jacob2018quantization}, knowledge distillation \cite{hinton2015distilling}, efficient architecture design \cite{tan2019efficientnet} and network pruning \cite{han2015learning,blalock2020state}. Pruning is attractive because it reduces parameters and computation of an existing architecture while preserving much of its accuracy \cite{liu2018rethinking}.

Classical pruning estimates the importance of each weight after dense training, either from the loss curvature, as in Optimal Brain Damage and Optimal Brain Surgeon \cite{lecun1989optimal,hassibi1992second}, or from weight magnitude \cite{han2015learning}, and then fine-tunes the pruned network \cite{liu2017learning}. Because this train--prune--fine-tune pipeline is expensive, \emph{sparse training} methods learn the sparse network during a single training run. Some fix the mask at initialisation \cite{lee2018snip,wang2020picking,tanaka2020synflow}, others prune gradually during training \cite{zhu2017prune,sanh2020movement}, and \emph{dynamic sparse training} (DST) methods keep the sparsity constant while periodically dropping and regrowing connections \cite{mocanu2018scalable,bellec2017deep,mostafa2019parameter,dettmers2019sparse,evci2020rigging}. In the most widely used DST method, RigL \cite{evci2020rigging}, connections are dropped by weight magnitude and regrown by gradient magnitude.

Magnitude says nothing about the signal a weight multiplies. A large weight attached to an input channel that is almost constant over the data contributes little beyond a bias, whereas a moderate weight on a channel with a wide, input-dependent distribution takes part in many different responses. Activation statistics have long been used to choose \emph{channels} to prune: the fraction of zero activations \cite{hu2016network}, the entropy of pooled channel responses \cite{luo2017entropy}, feature-map rank \cite{lin2020hrank} and channel independence \cite{sui2021chip}. We are not aware of work that evaluates them as the \emph{drop criterion of a DST loop}, where the score is recomputed many times during training and interacts with gradient-based regrowth. \TODO{Confirm this with a literature search (2022--2025 DST papers) before submission.}

This paper tests that idea. We propose Dynamic Entropy-Guided Sparse Training (DEGST), which estimates the histogram entropy $H(A)$ of each input channel from the training batches preceding a topology update, scores each weight by $H(A)\cdot|w|$, and otherwise follows a standard prune-and-regrow loop. Because the same entropy estimate is defined per channel, it also yields a channel score. We use it in DEGST-S, which prunes and regrows coupled channel groups and removes them physically, and in a hybrid variant that splits one weight budget between channel removal and unstructured sparsity.

A drop criterion that combines entropy with magnitude can only be credited to entropy if it beats magnitude alone under otherwise identical conditions. We therefore design the evaluation around three controls. First, an ablation runs the same loop with $|w|$, $H(A)$ and $H(A)\cdot|w|$ as the only difference. Second, all methods, including RigL, are compared at the same weight budget, with one shared training recipe and three seeds. Third, accuracy is reported against FLOPs rather than against nominal sparsity. The third control matters because methods that allocate sparsity globally prune late, low-resolution layers hardest: at equal parameter sparsity they remove far fewer FLOPs than layer-uniform methods, and a comparison at equal sparsity would mix two different operating regimes. We also report executed FLOPs and measured latency separately from non-zero FLOPs, because unstructured sparsity does not accelerate dense GPU kernels \cite{gale2020sparse,elsen2020fast,hoefler2021sparsity}.

The contributions of this paper are:
\begin{itemize}
    \item DEGST, a dynamic sparse training method whose drop criterion weights magnitude by the entropy of the input activation, with channel-level (DEGST-S) and hybrid variants that share the entropy estimate and remove channels physically using dependency-graph pruning \cite{fang2023depgraph}.
    \item A controlled ablation that separates the contribution of entropy from that of the prune-and-regrow loop, together with a comparison of further criteria (Taylor, Fisher, gradient, movement) inside the same loop.
    \item An evaluation at matched weight budgets on CIFAR-10 with ResNet-18, complemented by CIFAR-100 and MobileNetV2. It reports accuracy--FLOPs Pareto fronts for both non-zero and executed FLOPs, a documented latency/throughput protocol, training cost and hyper-parameter sensitivity.
    \item \TODO{One sentence summarising the main empirical finding, stated plainly (positive or negative), once C0 and the Pareto sweep are complete.}
\end{itemize}
